\section{Workload Characterization}
\label{sec:workloads}
\subsection{Human Activity Recognition}
The UCI Human Activity Recognition Using Smartphones dataset
\cite{uciHAR} supplies 128-sample windows with nine inertial channels and six
activity labels. Four-sample patches produce 32 tokens, each embedded from 36
values to width 128. The study uses sequential sensor channels rather than the
precomputed feature table. The recorded split contains 5,991 training, 1,361
validation, and 2,947 official test windows. Validation subjects are held out
from the official training subjects, and normalization is fitted on training
data only.

\subsection{Speech Commands}
The secondary workload uses Speech Commands v0.02 with all 35 named words
\cite{arxiv180403209}. Official lists define 84,843 training, 9,981 validation,
and 11,005 test clips. The host frontend converts mono 16 kHz audio into 40
log-mel channels in 64 frame slots. Two-frame patches again yield 32 tokens.
Two trailing log-silence frames are valid inputs in the frozen frontend; they
are not hardware padding. This is a 35-class protocol, not a directly
comparable 12-class keyword benchmark.

\subsection{Operator and Memory Demand}
For token count $N$, width $d$, and feed-forward expansion $r$, a block executes
\begin{equation}
  M_{\mathrm{block}} = 4Nd^2 + 2N^2d + 2rNd^2
\end{equation}
projection, attention, and feed-forward multiply--accumulates. These counts do
not include normalization, nonlinear functions, selection, division, residual
addition, or data movement. All effective operation counts must specify
whether summary construction is included.
\writingtask{architecture lead}{Add measured activation ranges, operator shares,
roofline and bandwidth analyses, and physical memory costs. Keep analytical
sequence sweeps distinct from hardware configurations actually implemented.}
